% Source: https://artofproblemsolving.com/community/c3h2694801_trapezoids
% Stable identity: aops-topic-2694801-q01
\emph{Height.} Through $D$ draw the line parallel to $BC$, meeting $AB$ at $E$. Then $DEBC$ is a parallelogram, so $DE=8$, $EB=5$ and $AE=5$. Triangle $ADE$ has $AD=AE=5$ and $DE=8$, so its altitude from $A$ is $\sqrt{25-16}=3$ and $[ADE]=\dfrac12\cdot8\cdot3=12$. Its altitude from $D$ is the height $h$ of the trapezoid: $\dfrac12\cdot5\cdot h=12$ gives $h=\dfrac{24}{5}$. Hence $[BCD]=\dfrac12\cdot CD\cdot h=\dfrac12\cdot5\cdot\dfrac{24}{5}=12$.

\emph{Ratio $BP:PD$.} Extend $AM$ to meet line $DC$ at $N$. Since $AB\parallel CN$ and $BM=MC$, triangles $MAB$ and $MNC$ are congruent, so $CN=AB=10$ and $DN=15$. Triangles $PAB$ and $PND$ are similar (as $AB\parallel DN$), so
\[\dfrac{BP}{PD}=\dfrac{AB}{DN}=\dfrac{10}{15}=\dfrac23,\qquad\dfrac{BP}{BD}=\dfrac25.\]

\emph{Area.} Triangles $BMP$ and $BCD$ share the angle at $B$, so
\[[BMP]=\dfrac{BM}{BC}\cdot\dfrac{BP}{BD}\cdot[BCD]=\dfrac12\cdot\dfrac25\cdot12=\boxed{\dfrac{12}{5}}.\]

\begin{asy}
usepackage("amsmath");
size(6.5cm);
defaultpen(fontsize(10pt));
pair A=(0,0),B=(10,0),C=(3.6,4.8),D=(-1.4,4.8),M=(B+C)/2,Nn=(13.6,4.8);
pair P=extension(A,M,B,D);
fill(B--M--P--cycle,lightgray);
draw(A--B--C--D--cycle,linewidth(0.8));draw(B--D);draw(A--M);draw(M--Nn,dashed);draw(C--Nn,dashed);
dot(M);dot(P);dot(Nn);
label("$A$",A,SW);label("$B$",B,SE);label("$C$",C,N);label("$D$",D,NW);label("$M$",M,E);label("$P$",P,N);label("$N$",Nn,NE);
label("$10$",(A+B)/2,S);label("$5$",(C+D)/2,N);label("$5$",(A+D)/2,W);label("$8$",B+0.3*(C-B),NE);
\end{asy}
